\documentclass[11pt]{article}
\usepackage[margin=1in]{geometry}
\usepackage{amsmath,amssymb,amsthm,booktabs,needspace}
\usepackage[hidelinks]{hyperref}
\setlength{\emergencystretch}{2em}
\newtheorem{theorem}{Theorem}
\newtheorem{proposition}[theorem]{Proposition}
\newcommand{\rank}{\operatorname{rank}}
\newcommand{\mr}{\operatorname{mr}}
\title{Two common double covers for eight retained MRC targets}
\author{Internal research working proof}
\date{17 November 2026}
\begin{document}
\maketitle

The preceding thirteen-sheet homology calculation suggested two four-knot
blocks. Exact graph certificates now prove Conway mutation inside each
block. Consequently these eight ordinary double-cover rank questions
reduce to two. Neither rank is determined here, and no additional knot
equality is proved. The finite inputs have an independent replay; the
topological implications are explicit classical imports. This working
proof is unformalized and externally unreviewed, with no novelty claim.

\begin{theorem}\label{thm:blocks}
Let
\[
 \mathcal A=\{13a_{129},13a_{136},13a_{194},13a_{197}\},\qquad
 \mathcal B=\{13a_{135},13a_{137},13a_{196},13a_{198}\}.
\]
Within each displayed block, the retained reduced alternating diagrams
are related by Conway mutations. In particular their double branched
covers are homeomorphic. Write the two fundamental groups as $G_A,G_B$.
Then
\[
 H_1(G_A;\mathbb Z)=C_{391},\qquad H_1(G_B;\mathbb Z)=C_{401},
 \qquad 2\le\rank G_A,\rank G_B\le3.
\]
If $\rank G_A=3$, every knot in $\mathcal A$ has $\mr(K)=b(K)=4$.
The same conditional conclusion holds for $G_B$ and $\mathcal B$.
\end{theorem}

\section{Classical inputs and their exact scope}

Greene's Corollary~4.5 states that connected reduced alternating
diagrams are mutants precisely when their Tait graphs are
2-isomorphic, with the fixed checkerboard convention of that paper.
Section~4 obtains this through graph operations and changes of plane
drawing. Parallel edges are permitted. We import that implication,
including its classical Whitney and plane-drawing inputs; the paper's
Floer and lattice classification is not needed. The read source is
the author preprint, printed pages~19--23; the published metadata is
\emph{Inventiones mathematicae}~192 (2013), 717--750,
\url{https://doi.org/10.1007/s00222-012-0421-4}.

The double-cover implication is Viro's classical observation, with
the complete short proof read in Morton's \emph{Mutant knots},
16~March~2014, Theorem~3, printed page~6. The two tangle covers
meet along a torus. A chosen lift of the Conway halfturn is isotopic
to the identity there, so mutation does not change the glued cover.
Viro's original publication receives credit; its full original proof
has not been audited. Morton's author manuscript is
\url{https://www.liverpool.ac.uk/~su14/papers/mutantchapter.pdf}.

These imports establish a homeomorphism and hence a group isomorphism.
No claim about mutation invariance of bridge number or meridional rank
is used. The machine certificate below is an abstract graph certificate,
not a literal sequence of modified planar diagrams or Tietze maps.
The classical graph-to-diagram theorem supplies that existence step.

\section{Frozen diagrams and finite certificates}

The frozen plan selects exactly the eight names in
Theorem~\ref{thm:blocks}; no source is replaced. Literal PD codes come
from the already pinned \texttt{parabolic\_panel.json}. Each has thirteen
crossings. The producer constructs both checkerboard graphs and selects
the smaller, with seven vertices and thirteen edges. Every source is a
single alternating knot on a sphere, and the selected graph is connected
with neither loops nor bridges. The other checkerboard graph is its
plane dual, so these checks imply reducedness.

Edges are numbered $0,\ldots,12$ in PD crossing order. For reproducibility,
the selected root graphs have the following edge lists:
\begin{align*}
 E_A={}&(01,12,20,13,31,04,43,35,56,06,36,36,50),\\
 E_B={}&(01,02,02,31,32,42,42,15,56,46,45,45,61).
\end{align*}
Here $uv$ denotes an undirected edge, and repeated entries are distinct
parallel edges. Source $13a_{129}$ is the root for $\mathcal A$ and
$13a_{135}$ for $\mathcal B$. Vertex names have no intrinsic meaning.

A Whitney switch uses an edge partition whose two sides meet in exactly
two vertices $u,v$, and interchanges $u,v$ on one side. The certificate
retains every boundary and side, followed by vertex and crossing-edge
bijections to the target graph. The search is bounded to $5000$ states
and thirty seconds per comparison. All six comparisons succeed:
\begin{center}
\begin{tabular}{lllrr}
\toprule
Block & Root & Target & Switches & States visited\\
\midrule
$\mathcal A$ & $13a_{129}$ & $13a_{136}$ &1&2\\
 & & $13a_{194}$ &1&2\\
 & & $13a_{197}$ &0&1\\
$\mathcal B$ & $13a_{135}$ & $13a_{137}$ &1&4\\
 & & $13a_{196}$ &0&1\\
 & & $13a_{198}$ &1&4\\
\bottomrule
\end{tabular}
\end{center}
Zero switches means that the abstract graphs are isomorphic. It does
not assert that their plane embeddings or knot diagrams coincide.

The local crossing convention is tracked explicitly. In the producer's
face-dart coordinates, the selected opposite corners have parity zero
throughout block $\mathcal A$ and parity one throughout block $\mathcal B$.
Thus each comparison uses the same convention at every crossing.
If that convention is the dual of Greene's choice, dualizing both
graphs preserves their cycle-matroid isomorphism, since matroid duality
preserves isomorphisms. No target needs mirroring. The independent corner
construction below also checks this parity rather than assuming it.

\section{Independent replay}

\texttt{verify\_mutation\_covers.py} imports none of the producer's graph,
canonicalization or search routines. It identifies regions by gluing
corners across diagram arcs, rather than by the producer's face
permutation. It checks the single alternating knot walk, Euler region
count, checkerboard bipartition and source PD identity. Complete vertex
stars align its independently constructed graph with the retained one.
Enumeration of all spanning trees verifies connectedness and the
absence of loops and bridges, giving $391$ trees for every graph in
$\mathcal A$ and $401$ in $\mathcal B$.

For a graph with edges $e_i=u_iv_i$, it enumerates all $2^{13}=8192$
edge subsets and retains exactly those with even incidence at every
vertex. These form its full binary cycle space, with
$2^{13-7+1}=128$ elements for every named source. Every recorded switch
has the claimed two-vertex boundary and preserves that space with edge
labels fixed. The final edge bijection transports the entire space
onto the target's, and the final vertex bijection is checked on every
edge. Inclusion-minimal nonzero vectors are the graph circuits, so this
proves the required cycle-matroid isomorphism, not just equality of
polynomials or tree counts.

Two positive controls replay: an explicit thirteen-edge graph with a
genuine switch changing its degree multiset, and literal pretzel
diagrams $P(2,3,3,5)$ and $P(2,3,5,3)$. Four corruptions are rejected:
a wrong switch boundary, a bijective but wrong crossing transport, a
wrong final vertex map, and a single crossing change breaking the
alternating convention. The retained receipt reports eight named
diagrams, six comparisons, two controls and four rejected corruptions.

The finite certificate is in \texttt{mutation\_cover\_results.json};
the independent receipt is \texttt{mutation\_cover\_verification.json}.
\texttt{mutation\_source\_audit.json} seals the two read classical
sources. The separate archived KnotInfo subset records equal Jones
and Conway polynomials and signatures within the blocks. Those
equalities motivated screening; they are not the mutation proof.

\section{Rank consequence and remaining question}

\begin{proof}[Proof of Theorem~\ref{thm:blocks}]
The independently checked graph certificates, reduced alternating
diagram checks and common crossing conventions satisfy Greene's
Corollary~4.5. Hence the six root-to-member diagram pairs are mutants.
Morton's Theorem~3, credited to Viro, identifies their double branched
covers. The retained integral homology calculation gives the displayed
cyclic groups, so the two blocks cannot have isomorphic cover groups.

Each root has a retained epimorphism onto the nonabelian group
$\mathrm{SL}_3(\mathbb F_3)$, proving rank at least two. The already
certified four-meridian upper and the double-cover inequality give
\[
 \rank\pi_1(\Sigma_2(K))\le\mr(K)-1\le b(K)-1\le3.
\]
If either common rank is three, the same inequality forces
$4\le\mr(K)\le b(K)\le4$ for every knot of its block.
Each upper bound is an individual retained certificate; it is not
transferred by an assumed mutation-invariance theorem.
\end{proof}

The rank-three obstruction is supplied in
\path{FOUR_PIECE_DOUBLE_COVERS_SHARP_FOUR.tex}, using the compact
presentations and certified peripheral models. Nonabelianness alone
proves only rank at least two; the three-sheet cover criterion is a
different numerical gate. The three-piece and final-five companions
settle the other original eight targets. The complete sixteen-target
result and its existing graph-cover family scope are recorded in
\path{ORIGINAL_SIXTEEN_KNOWN_GRAPH_COVER_FAMILY.tex}. Negative
quotient-module tests do not imply rank two, and full MRC remains open.

\end{document}
